% ----------------------------------------------------------
\begin{frame}{Constraints -- Overview}
  Constraints are \textbf{rules enforced by the DBMS} that prevent invalid data from entering the
  database. They are the last line of defence for data integrity -- independent of the application.

  \begin{center}
  \renewcommand{\arraystretch}{1.3}
  \begin{tabular}{@{}lll@{}}
    \toprule
    \textbf{Constraint}  & \textbf{Scope}  & \textbf{Guarantees} \\
    \midrule
    \texttt{NOT NULL}    & Column & Column always contains a value \\
    \texttt{UNIQUE}      & Column/Table & All values are distinct \\
    \texttt{PRIMARY KEY} & Column/Table & NOT NULL + UNIQUE; identifies each row \\
    \texttt{FOREIGN KEY} & Column & References a valid row in another table \\
    \texttt{CHECK}       & Column/Table & Boolean expression is always true \\
    \texttt{DEFAULT}     & Column & Fallback value when none is provided \\
    \bottomrule
  \end{tabular}
  \end{center}

  \begin{block}{Why Not Just Validate in Application Code?}
    Application code can be bypassed (direct DB access, bugs, multiple applications).
    Constraints enforced at the \textbf{database level} hold regardless of which client writes data.
  \end{block}
\end{frame}

% ----------------------------------------------------------
\begin{frame}[fragile]{Constraints in Practice -- Full Example}
  \begin{lstlisting}[style=sqlstyle]
CREATE TABLE Track (
    id           INT            AUTO_INCREMENT,
    title        VARCHAR(200)   NOT NULL,
    duration_sec INT            NOT NULL DEFAULT 0,
    track_no     TINYINT        UNSIGNED,
    album_id     INT            NOT NULL,
    CONSTRAINT pk_track
        PRIMARY KEY (id),
    CONSTRAINT uq_track_pos
        UNIQUE (album_id, track_no),       -- no duplicate positions
    CONSTRAINT fk_track_album
        FOREIGN KEY (album_id) REFERENCES Album(id)
        ON DELETE CASCADE ON UPDATE CASCADE,
    CONSTRAINT chk_duration
        CHECK (duration_sec >= 0)
);
  \end{lstlisting}

  \begin{itemize}
    \item \texttt{UNIQUE (album\_id, track\_no)}: composite unique -- each track number
          must be unique \emph{within} an album, but the same number can occur in different albums.
  \end{itemize}
\end{frame}

% ----------------------------------------------------------
\begin{frame}[fragile]{Foreign Keys -- ON DELETE Actions}
  What should happen to child rows when the parent row is deleted?

  \begin{lstlisting}[style=sqlstyle]
-- CASCADE: deleting Artist also deletes all their Albums
FOREIGN KEY (artist_id) REFERENCES Artist(id)
    ON DELETE CASCADE

-- SET NULL: deleting Artist sets album.artist_id to NULL
FOREIGN KEY (artist_id) REFERENCES Artist(id)
    ON DELETE SET NULL

-- RESTRICT (default): reject the DELETE if child rows exist
FOREIGN KEY (artist_id) REFERENCES Artist(id)
    ON DELETE RESTRICT

-- SET DEFAULT: sets child FK to its DEFAULT value
FOREIGN KEY (artist_id) REFERENCES Artist(id)
    ON DELETE SET DEFAULT
  \end{lstlisting}
\end{frame}

% ----------------------------------------------------------
\begin{frame}{ON DELETE Comparison Table}
  \begin{center}
  \renewcommand{\arraystretch}{1.4}
  \begin{tabularx}{\linewidth}{@{}lXX@{}}
    \toprule
    \textbf{Action} & \textbf{Effect on child rows} & \textbf{Use when \ldots} \\
    \midrule
    \texttt{RESTRICT}   & Delete is \textbf{blocked} if child rows exist & Child rows must not become orphans (default) \\[4pt]
    \texttt{CASCADE}    & Child rows are \textbf{deleted automatically} & Children are meaningless without the parent (e.g.\ order lines without an order) \\[4pt]
    \texttt{SET NULL}   & FK column set to \textbf{NULL}                & Children can exist independently (e.g.\ products after a brand is discontinued) \\[4pt]
    \texttt{SET DEFAULT}& FK column set to its \textbf{DEFAULT value}   & A fallback parent exists (rare) \\
    \bottomrule
  \end{tabularx}
  \end{center}

  \begin{alertblock}{CASCADE -- Handle with Care}
    \texttt{ON DELETE CASCADE} propagates deletes automatically and silently. Deleting one top-level
    row can cascade through multiple related tables. Always verify which tables are affected before
    using CASCADE in a production schema.
  \end{alertblock}
\end{frame}

% ----------------------------------------------------------
\begin{frame}[fragile]{ON DELETE -- Concrete Examples}
  \begin{lstlisting}[style=sqlstyle]
-- Scenario 1: CASCADE
-- Deleting artist "Aurora" also deletes all her Albums and
-- (if Track has CASCADE too) all their Tracks.
DELETE FROM Artist WHERE name = 'Aurora';
-- -> Album rows with artist_id = Aurora.id are deleted
-- -> Track rows of those albums are deleted too

-- Scenario 2: SET NULL
-- Artist is deleted; albums remain but have no artist
DELETE FROM Artist WHERE id = 1;
-- -> Album.artist_id becomes NULL for those albums

-- Scenario 3: RESTRICT (default)
-- Cannot delete artist if albums reference them
DELETE FROM Artist WHERE id = 1;
-- ERROR: Cannot delete or update a parent row:
-- a foreign key constraint fails (Album.fk_album_artist)
  \end{lstlisting}
\end{frame}

% ----------------------------------------------------------
\begin{frame}[fragile]{Named Constraints and ALTER TABLE ADD CONSTRAINT}
  \begin{lstlisting}[style=sqlstyle]
-- Naming constraints at table creation
CREATE TABLE Album (
    id        INT,
    title     VARCHAR(200),
    year      YEAR,
    artist_id INT,
    CONSTRAINT pk_album    PRIMARY KEY (id),
    CONSTRAINT chk_year    CHECK (year BETWEEN 1877 AND 2100)
);

-- Add a constraint to an existing table
ALTER TABLE Album
ADD CONSTRAINT fk_album_artist
    FOREIGN KEY (artist_id) REFERENCES Artist(id)
    ON DELETE CASCADE;

-- Remove a named constraint
ALTER TABLE Album DROP FOREIGN KEY fk_album_artist;
ALTER TABLE Album DROP CHECK chk_year;
  \end{lstlisting}

  \begin{block}{Why Name Constraints?}
    Named constraints appear in error messages by name (e.g.\ ``Duplicate entry violates
    \texttt{uq\_track\_pos}''), making debugging faster. They can also be dropped individually.
  \end{block}
\end{frame}
